
\section{Proof of the Convergence of AGMs}\label{app:proof-convergence}

% \xinghan{introduce mu, and assume gamma / beta1}

In this section, we aim to prove \Cref{thm:formal-convergence-1} and the first part of \Cref{thm:formal-main-result}. Specifically, for some constant $\gamma = 1 - \Theta(\eta)$, we prove that the loss value of AGM converges to $\tilde{\O}(\gamma^K + \eta)$ within $K$ steps with high probability. If we substitute $K = \O\left(\frac{1}{\eta}\log \frac{1}{\eta}\right)$, this will recover the first part of \Cref{thm:formal-main-result}; However, this convergence analysis works for any $K = \O(\mathrm{poly}(1/\eta))$, and substituting $K = \O(\eta^{-2})$ will give us a high probability guarantee that the iteration stays near manifold in the whole scope of our analysis, which helps the proof of the second part too.

First, we introduce some additional notations that will be used in our proof. In the AGM framework, an algorithm starts from initial state $\vtheta_0$, and we set $\vm_0 = \bf{0} \in \R^d$, $\vv_0 = \bf{0} \in \R^D$. For every $k \geq 0$, we use \textbf{step $k+1$} to refer to the process of obtaining the noisy gradient $\nabla \ell_k(\vtheta_{k})$ and then $\vm_{k+1}, \vv_{k+1}$ and $\vtheta_{k+1}$. 

For any $k \geq 0$, to simplify the notation, we denote that
\begin{gather*}
    \vg_{k} := \nabla \ell_k(\vtheta_{k}), \quad
    \vz_{k} := \ell_k(\vtheta_{k}) - \cL(\vtheta_k) \sim \cZ(\vtheta_k), \quad \mS_k :=S(\vv_{k}),\\
    \mU_{k+1} := S(\vv_{k}) \vg_{k}, \quad
    \vu_{k+1} := S(\vv_{k}) \vm_{k+1},\quad \vphi_k:= \Phi_{\mS_k}(\vtheta_k).
\end{gather*} 
We use \textbf{time $k$} to refer to the time right before step $k + 1$ happens, i.e. the time right after we get $\vtheta_k$. We also define $\left\{\mathcal{F}_k\right\}$ as the natural filteration generated by the history of optimization, where each $\mathcal{F}_k = \sigma\left(\vtheta_0, \vz_0, \cdots, \vz_{k-1}\right)$ can be interpreted as ``all the information available up to time $k$''. We use the notation \textbf{$\E_k$} to denote the expectation conditioned on $\mathcal{F}_k$. 

To start with, we prove that the descent direction of each step does not veer off the direction of a preconditioned gradient descent, and the mismatch term can also be constrained by a list of martingales. After that, we can ensure a decay in the loss function every step, with some small perturbations that can be dealt with using Azuma-Hoeffding's inequality.

From \cref{descent-direction} throughout \cref{convergence}, we will assume that the loss function  $\cL$ satisfies $\mu$-PL condition at each iteration step, which is automatically satisfied in the setting of \Cref{thm:formal-convergence-1}. follows directly from the result. After that, we argue that if the loss function satisfies $\mu$-PL only within some local neighborhood, an AGM starting near enough to the manifold will stick to the manifold with high probability, which leads to the first part of \Cref{thm:formal-main-result}. 
% \inp{}{}
\begin{lemma}\label{descent-direction}
    Let $\cL$ satisfy \cref{amp:smooth}. Define $\tilde{\vv}_{k} := \beta_2 {\vv}_{k-1} + (1-\beta_2) \E_{k-1}[V\left(\vg_{k-1}\vg_{k-1}^\top\right)]$. There exist a constant $C_{1a}$ and a constant $C_{1b}$ independent of $\cL$, such that for any $k \geq 1$,
  \begin{align*}
\left\langle\nabla\cL\left(\vtheta_{k-1}\right), \mU_{k}\right\rangle = \nabla\cL\left(\vtheta_{k-1}\right)^\top S(\tilde{\vv}_k) \nabla\cL\left(\vtheta_{k-1}\right) - Y_{k} - X_{k},
\end{align*}
where $Y_k$ and $X_k$ are two $\mathcal{F}_{k}$-measurable random variables such that: \begin{enumerate}
    \item $\left|Y_k\right| \leq C_{1a} \lnormtwo{\nabla \cL(\vtheta_{k-1})} \cdot \eta^2$ a.s.
    \item $\left|X_k\right| \leq C_{1b}\lnormtwo{\nabla \cL(\vtheta_{k-1})}$ a.s., and $ \E_{k-1}[X_k] = 0$.
\end{enumerate}
\end{lemma}
\begin{proof}
    We first peel the $S(\tilde{\vv}_k)$ part off the $S(\vv_k)$ term: 
    \begin{align*}
        \left\langle\nabla\cL\left(\vtheta_{k-1}\right), \mU_k\right\rangle &= \left\langle\nabla\cL\left(\vtheta_{k-1}\right), S(\vv_k) \vg_{k-1}\right\rangle \\
        &= \left\langle\nabla\cL\left(\vtheta_{k-1}\right), S(\tilde{\vv}_k) \vg_{k-1}\right\rangle + \left\langle\nabla\cL\left(\vtheta_{k-1}\right), \left(S(\vv_k) - S(\tilde{\vv}_k)\right) \vg_{k-1}\right\rangle.
    \end{align*}
    Define $Y_k$ as $Y_k = -\left\langle\nabla\cL\left(\vtheta_{k-1}\right), \left(S(\vv_k) - S(\tilde{\vv}_k)\right) \vg_{k-1}\right\rangle$, then it holds almost surely that 
    $$\left|Y_k \right| \leq \lnormtwo{\nabla \cL(\vtheta_{k-1})}\lnormtwo{\left(S(\vv_k) - S(\tilde{\vv}_k)\right) \vg_{k-1}}.$$ 
    Since $S$ is Lipscitz, $V$ is linear and 
    $$\lnormtwo{\tilde{\vv}_k - \vv_k} = (1-\beta_2)\lnormtwo{\E_{k-1}\left[V\left(\vg_{k-1}\vg_{k-1}^\top\right)\right] - V\left(\vg_{k-1}\vg_{k-1}^\top\right)},$$ we conclude that $\left| Y_k \right| \leq C_{1a} \lnormtwo{\nabla \cL(\vtheta_{k-1})}\cdot \eta^2$ a.s. for some constant $C_{1a}$. The rest term $\left\langle\nabla\cL\left(\vtheta_{k-1}\right), S(\tilde{\vv}_k) \vg_{k-1}\right\rangle$ can also be decomposed into a deterministic part and a random part as: \begin{align*}
        \left\langle\nabla\cL\left(\vtheta_{k-1}\right), S(\tilde{\vv}_k) \vg_{k-1}\right\rangle &= \left\langle\nabla\cL\left(\vtheta_{k-1}\right), S(\tilde{\vv}_k) \left(\nabla \cL(\vtheta_{k-1}) + \vz_{k-1}\right)\right\rangle \\
        &= \nabla\cL\left(\vtheta_{k-1}\right)^\top S(\tilde{\vv}_k) \nabla\cL\left(\vtheta_{k-1}\right) + \left\langle \vz_{k-1}, S(\tilde{\vv}_k)^\top \nabla\cL\left(\vtheta_{k-1}\right)\right\rangle.
    \end{align*}
    Now we only need to let $X_k = \left\langle \vz_{k-1}, S(\tilde{\vv}_k)^\top \nabla\cL\left(\vtheta_{k-1}\right)\right\rangle$. It's easy to see that $\E_{k-1}[X_k] = 0$ and $\left| X_k \right| \leq C_{1b} \lnormtwo{\nabla \cL(\vtheta_{k-1})}$ a.s. for some constant $C_{1b}$. Finally, note that $C_{1b}$ is the multiplication of a constant bounding the magnitude of $\vz$ and $R_2$ which bounds $\normtwo{S}$, which is independent of $\cL$. This completes the proof.
\end{proof}

\begin{lemma}[Descent Lemma of the AGM Framework]\label{descent-lemma}
    Let $\cL$ satisfy \cref{amp:smooth}. For any $k \geq 1$ it holds that \begin{align*}
        \cL(\vtheta_{k}) - \cL(\vtheta_{k-1}) \leq C_2 \eta^2 -\eta (1-\beta_1)\sum_{i=1}^{k} \beta_1^{k-i}\left\langle \nabla\cL(\vtheta_{i-1}), \mU_i \right\rangle
    \end{align*}
    for some constant $C_2$.
\end{lemma}
\begin{proof}
    From the smoothness of $\cL$ we have \begin{align*}
        \cL(\vtheta_{k}) - \cL(\vtheta_{k-1}) & \leq - \left\langle \nabla \cL(\vtheta_{k-1}), \eta \vu_k \right\rangle + \frac{\rho\eta^2}{2} \lnormtwo{\vu_k}^2.
    \end{align*}   
    If $k = 1$, then $\vm_k$ = $(1-\beta_1)\vg_{k-1}$, so $\vu_k = (1-\beta_1)\mU_{k}$, and the statement trivially holds as long as $C_2 \geq \frac{\rho}{2}\lnormtwo{\vu_k}^2$. If $k > 1$, then the $- \left\langle \nabla \cL(\vtheta_{k-1}), \vu_k \right\rangle$ term can be expanded as \begin{align*}
        - \left\langle \nabla \cL(\vtheta_{k-1}), \vu_k \right\rangle &= - \left\langle \nabla \cL(\vtheta_{k-1}), S(\vv_k)\vm_k \right\rangle \\
        &=- \left\langle \nabla \cL(\vtheta_{k-1}), S(\vv_k)\left(\beta_1 \vm_{k-1} + (1-\beta_1)\vg_{k-1}\right) \right\rangle \\
        &=- \beta_1\left\langle \nabla \cL(\vtheta_{k-1}), S(\vv_k)\vm_{k-1}\right\rangle - (1-\beta_1)\left\langle \nabla \cL(\vtheta_{k-1}),S(\vv_k)\vg_{k-1}\right\rangle \\
        &=- \beta_1\left\langle \nabla \cL(\vtheta_{k-2}), S(\vv_{k-1})\vm_{k-1}\right\rangle - (1-\beta_1)\left\langle \nabla \cL(\vtheta_{k-1}),\mU_k\right\rangle \\
        &\quad - \beta_1 \left\langle \nabla\cL(\vtheta_{k-1})-\nabla\cL(\vtheta_{k-2}),S(\vv_{k-1})\vm_{k-1}\right\rangle \\
        &\quad - \beta_1 \left\langle \nabla\cL(\vtheta_{k-1}),\left(S(\vv_{k})-S(\vv_{k-1})\right)\vm_{k-1}\right\rangle \\
        &\leq - \beta_1\left\langle \nabla \cL(\vtheta_{k-2}), S(\vv_{k-1})\vm_{k-1}\right\rangle - (1-\beta_1)\left\langle \nabla \cL(\vtheta_{k-1}),\mU_k\right\rangle \\
        &\quad + \beta_1 \lnormtwo{\nabla\cL(\vtheta_{k-1})-\nabla\cL(\vtheta_{k-2})}\lnormtwo{S(\vv_{k-1})\vm_{k-1}} \\
        &\quad + \beta_1 \lnormtwo{\nabla\cL(\vtheta_{k-1})}\lnormtwo{\left(S(\vv_{k})-S(\vv_{k-1})\right)\vm_{k-1}}.
    \end{align*}
    Note that a single step of update on $\vtheta$ and $\vv$ is small since \begin{align*}
        \vtheta_{k} - \vtheta_{k-1} &= \eta \vu_k, \\
        \vv_{k} - \vv_{k-1} &= \beta_2 \vv_{k-1} + (1-\beta_2) V\left(\vg_{k-1} \vg_{k-1}^\top\right) - \vv_{k-1} \\
        &= (1-\beta_2) \left(V\left(\vg_{k-1} \vg_{k-1}^\top\right) - \vv_{k-1}\right),
    \end{align*}  
    which implies that $\lnormtwo{\vtheta_{k} - \vtheta_{k-1}} = \O(\eta)$ and $\lnormtwo{\vv_{k} - \vv_{k-1}} = \O(\eta^2)$. We then leverage the smoothness of $\nabla \cL$ and $S$ to conclude that there exists some constant $\tilde{C}_2$ such that
    \begin{align*}
    - \left\langle \nabla \cL(\vtheta_{k-1}), \vu_k \right\rangle 
        & \leq - \beta_1\left\langle \nabla \cL(\vtheta_{k-2}), \vu_{k-1}\right\rangle - (1-\beta_1)\left\langle \nabla \cL(\vtheta_{k-1}),\mU_k\right\rangle + \beta_1 \tilde{C}_2 \eta.
    \end{align*} 
    Giving that $\vu_{0} = \bold{0}$, we can expand this formula iteratively as \begin{align*}
        - \left\langle \nabla \cL(\vtheta_{k-1}), \vu_k \right\rangle 
        & \leq - \beta_1\left\langle \nabla \cL(\vtheta_{k-2}), \vu_{k-1}\right\rangle - (1-\beta_1)\left\langle \nabla \cL(\vtheta_{k-1}),\mU_k\right\rangle + \beta_1 \tilde{C}_2 \eta \\
        & \leq -\beta_1^2 \left\langle \nabla \cL(\vtheta_{k-3}), \vu_{k-2}\right\rangle - \beta_1(1-\beta_1)\left\langle \nabla \cL(\vtheta_{k-2}),\mU_{k-1}\right\rangle \\
        & \quad - (1-\beta_1)\left\langle \nabla \cL(\vtheta_{k-1}),\mU_k\right\rangle + \beta_1\tilde{C}_2 \eta + \beta_1^2\tilde{C}_2 \eta \\
        & \leq \cdots \\
        & \leq -(1-\beta_1)\sum_{i=1}^{k} \beta_1^{k-i}\left\langle \nabla\cL(\vtheta_{i-1}), \mU_i \right\rangle + \beta_1^{k-i+1} \tilde{C}_2 \eta \\
        & \leq \frac{\beta_1}{1-\beta_1} \tilde{C}_2 \eta - (1-\beta_1)\sum_{i=1}^{k} \beta_1^{k-i}\left\langle \nabla\cL(\vtheta_{i-1}), \mU_i \right\rangle.
    \end{align*}
    Plugging in, we get \begin{align*}
        \cL(\vtheta_{k}) - \cL(\vtheta_{k-1}) & \leq \frac{\beta_1}{1-\beta_1} \tilde{C}_2 \eta^2 + \frac{\rho\eta^2}{2} \lnormtwo{\vu_k}^2 -\eta (1-\beta_1)\sum_{i=1}^{k} \beta_1^{k-i}\left\langle \nabla\cL(\vtheta_{i-1}), \mU_i \right\rangle   \\
        &\leq  C_2 \eta^2 -\eta (1-\beta_1)\sum_{i=1}^{k} \beta_1^{k-i}\left\langle \nabla\cL(\vtheta_{i-1}), \mU_i \right\rangle,
    \end{align*}
    for some constant $C_2$.
\end{proof}


\begin{lemma}\label{descent-lemma-exp}
    Let $\cL$ satisfy \cref{amp:smooth}, and assume that $\mu$-PL condition is satisfied at all $\vtheta_k$ where $k \geq 0$. Define $\gamma: = 1 - \frac{2\eta\mu (1-\beta_1)}{R_0}$.
    For any $k \geq 0$, we have
    \begin{align*}
        \cL(\vtheta_{k})-\cL^* & \leq \gamma^k \left(\cL(\vtheta_0) - \cL^*\right) + \eta (1-\beta_1)\sum_{i=1}^{k}X_i  \sum_{j=i}^{k}\gamma^{k-j}\beta_1^{j-i} + C_{3} \eta
    \end{align*}
    for some constant $C_3$.
\end{lemma}
\begin{proof}
    We start from \cref{descent-lemma} and plug in \cref{descent-direction}:\begin{align*}
    \cL(\vtheta_{k}) - \cL(\vtheta_{k-1}) &\leq C_2 \eta^2 -\eta (1-\beta_1)\sum_{i=1}^{k} \beta_1^{k-i}\left\langle \nabla\cL(\vtheta_{i-1}), \mU_i \right\rangle \\
    & = C_2 \eta^2 - \eta (1-\beta_1) \sum_{i=1}^{k} \beta_1^{k-i} \left(\nabla\cL\left(\vtheta_{i-1}\right)^\top S(\tilde{\vv}_i) \nabla\cL\left(\vtheta_{i-1}\right) - Y_i - X_i\right).
    \end{align*}
    Since $\left|Y_i\right| \leq C_{1a}\lnormtwo{\nabla\cL(\vtheta_{i-1})}\cdot \eta^2$ for every $i$, the effect of $Y$ is negligible: \begin{align*}
        \left|\eta(1-\beta_1)\sum_{i=1}^{k}\beta_1^{k-i}Y_i\right| \leq C_{1a} \eta^3\cdot \max_{i=0}^{k}\left\{\lnormtwo{\nabla\cL(\vtheta_{i-1})}\right\} = o(\eta^2),
    \end{align*}
    and we can absorb it into the $C_2 \eta^2$ term to write out that 
    \begin{align*}
    \cL(\vtheta_{k}) - \cL(\vtheta_{k-1}) \leq \tilde{C}_3 \eta^2 - \eta (1-\beta_1) \sum_{i=1}^{k} \beta_1^{k-i} \left(\nabla\cL\left(\vtheta_{i-1}\right)^\top S(\tilde{\vv}_i) \nabla\cL\left(\vtheta_{i-1}\right) - X_i\right)
    \end{align*}
    for some constant $\tilde{C}_3$.
    Note that $S(\tilde{\vv}_i) \succeq \frac{1}{R_0}$, so $\nabla\cL\left(\vtheta_{i-1}\right)^\top S(\tilde{\vv}_i) \nabla\cL\left(\vtheta_{i-1}\right) \geq \frac{1}{R_0}\lnormtwo{\nabla\cL\left(\vtheta_{i-1}\right)}^2$ for any $i$, hence\begin{align}\label{equ:loss-decrease}
        \cL(\vtheta_{k}) - \cL(\vtheta_{k-1}) \leq \tilde{C}_3 \eta^2 - \eta (1-\beta_1) \sum_{i=1}^{k} \beta_1^{k-i} \left(\frac{1}{R_0}\lnormtwo{\nabla\cL\left(\vtheta_{i-1}\right)}^2 - X_i\right).
    \end{align}
    Combining with the $\mu$-PL property $\lnormtwo{\nabla\cL\left(\vtheta_{i-1}\right)}^2 \geq 2\mu\left(\cL(\vtheta_{i-1}) - \cL^* \right)$, we have
    % $$\cL(\vtheta_{k}) - \cL(\vtheta_{k-1}) \leq  C_2 \eta^2 - \frac{2\eta\mu (1-\beta_1)}{R_0} \sum_{i=1}^{k} \beta_1^{k-i} \left(\cL(\vtheta_{i-1}) - \cL^* \right) + \eta(1-\beta_1)\sum_{i=1}^{k}\beta_1^{k-i}(Y_i + X_i).$$
    \begin{align*}
        \cL(\vtheta_{k}) - \cL^* &\leq  \tilde{C}_3 \eta^2 + \cL(\vtheta_{k-1}) - \cL^* - \frac{2\eta\mu (1-\beta_1)}{R_0} \sum_{i=1}^{k} \beta_1^{k-i} \left(\cL(\vtheta_{i-1}) - \cL^* \right) \\
        & \quad + \eta(1-\beta_1)\sum_{i=1}^{k}\beta_1^{k-i}X_i \\
        &\leq \tilde{C}_3 \eta^2 + \left(1 - \frac{2\eta\mu (1-\beta_1)}{R_0}\right) \left(\cL(\vtheta_{k-1})-\cL^*\right) + \eta(1-\beta_1)\sum_{i=1}^{k}\beta_1^{k-i}X_i \\
        & = \tilde{C}_3 \eta^2 + \gamma \left(\cL(\vtheta_{k-1})-\cL^*\right) + \eta(1-\beta_1)\sum_{i=1}^{k}\beta_1^{k-i}X_i. 
    \end{align*}
    Note that we can expand the $\cL(\vtheta_{k-1})-\cL^*$ term iteratively to obtain a generic formula for $\cL(\vtheta_k) - \cL^*$:\begin{align*}
        \cL(\vtheta_k) - \cL^* & \leq \gamma \left(\cL(\vtheta_{k-1}) - \cL^*\right) + \eta(1-\beta_1)\sum_{i=1}^{k}\beta_1^{k-i}X_i + \tilde{C}_3 \eta^2 \\
        & \leq \gamma^k \left(\cL(\vtheta_0) - \cL^*\right) + \eta (1-\beta_1)\sum_{j=1}^{k}\gamma^{k-j}\sum_{i=1}^{j}\beta_1^{j-i}X_i + \sum_{j=1}^{k}\gamma^{k-j}\tilde{C}_3 \eta^2 \\
        &  \leq \gamma^k \left(\cL(\vtheta_0) - \cL^*\right) + \eta (1-\beta_1)\sum_{i=1}^{k}X_i  \sum_{j=i}^{k}\gamma^{k-j}\beta_1^{j-i} + C_{3} \eta,
    \end{align*}
    where $C_{3} = \tilde{C}_3 \cdot \frac{R_0}{2\mu(1-\beta_1)}$.
\end{proof}

\begin{corollary}\label{cor:nablaL-large}
    Let $\cL$ satisfy \cref{amp:smooth}. There exists a constant $\bar{C}$ independent of $\cL$, such that $\forall k > 0$, if \begin{align*}
        \normtwo{\nabla \cL(\vtheta_{k-1})} > \bar{C},
    \end{align*}
    then with sufficiently small $\eta$, we have \begin{align*}
        \cL(\vtheta_k) < \cL(\vtheta_{k-1}).
    \end{align*}
\end{corollary}
\begin{proof}
    Note that \Cref{equ:loss-decrease} can be obtained without PL condition: \begin{align*}
        \cL(\vtheta_{k}) - \cL(\vtheta_{k-1}) \leq \tilde{C}_3 \eta^2 - \eta (1-\beta_1) \sum_{i=1}^{k} \beta_1^{k-i} \left(\frac{1}{R_0}\lnormtwo{\nabla\cL\left(\vtheta_{i-1}\right)}^2 - X_i\right).
    \end{align*}
    From \cref{descent-direction}, $\left|X_i\right| \leq C_{1b}\lnormtwo{\nabla \cL(\vtheta_{i-1})}, \forall i \leq k$ where $C_{1b}$ is a constant independent of $\cL$. So
\begin{align*}
\frac{1}{R_0}\lnormtwo{\nabla\cL(\vtheta_{i-1})}^2 - X_i
\;\ge\; \frac{1}{R_0}\lnormtwo{\nabla\cL(\vtheta_{i-1})}^2 - C_{1b}\lnormtwo{\nabla\cL(\vtheta_{i-1})}
\;\ge\; -\,\frac{R_0 C_{1b}^2}{4},
\end{align*}
where the last inequality uses $a^2 - 2ab \ge -b^2$ with $a=\lnormtwo{\nabla\cL(\vtheta_{i-1})}/\sqrt{R_0}$ and
$b=C_{1b}\sqrt{R_0}/2$. 
Let $G_{i-1}:=\lnormtwo{\nabla\cL(\vtheta_{i-1})}$, then
\begin{align*}
\sum_{i=1}^{k} \beta_1^{k-i}\left(\frac{1}{R_0}G_{i-1}^2 - X_i\right)
&\ge \left(\frac{1}{R_0}G_{k-1}^2 - C_{1b} G_{k-1}\right)
   - \frac{R_0 C_{1b}^2}{4}\sum_{i=1}^{k-1} \beta_1^{k-i} \\
&\ge \frac{1}{R_0}G_{k-1}^2 - C_{1b} G_{k-1}
   - \frac{R_0 C_{1b}^2}{4}\cdot\frac{\beta_1}{1-\beta_1}.
\end{align*}
As long as $\frac{1}{R_0}G_{k-1}^2 - C_{1b} G_{k-1}
   - \frac{R_0 C_{1b}^2}{4}\cdot\frac{\beta_1}{1-\beta_1} > 0$, a small $\eta$ can ensure the loss strictly decreases at this step $k$. Set 
\begin{align*}
    \bar{C} := \frac{R_0 C_{1b}}{2}\left(1 + \sqrt{\frac{1}{1-\beta_1}}\right),
\end{align*}
then any $G_{k-1} > \bar{C}$ meets this requirement. Moreover, $\bar{C}$ depends only on $(R_0,C_{1b},\beta_1)$ and is independent of $\cL$. This proves the corollary.
\end{proof}

\begin{lemma}\label{martingale-bound-by-L}Let $\cL$ satisfy \cref{amp:smooth}, and assume that $\mu$-PL condition is satisfied at all $\vtheta_k$ where $k \geq 0$. Let $k \leq K = \O(\mathrm{poly}(1/\eta))$
     % \O\left(\frac{1}{\eta}\log\frac{1}{\eta}\right)$, 
     and let \begin{align*}\psi: \left(\{0,1,\cdots,k-1\}\times (0,1)\right) \longrightarrow \R^+\end{align*} be a function. Let $\left\{X_i\right\}_{i=1}^{k}$ be any martingale difference sequence such that: \begin{enumerate}
         \item $X_i$ is $\mathcal{F}_{i}$-measurable and $\E_{i-1}[X_i] = 0$;
         \item $\left|X_i\right| \leq C_{1b}\lnormtwo{\nabla\cL(\vtheta_{i-1})}$ a.s.
     \end{enumerate} for any $i \in [k]$. If for any $i \in [k]$ and $\delta \in (0, 1)$, it holds with probability $1-\delta$ that $$\cL(\vtheta_{i-1})-\cL^* \leq \psi(i, \delta),$$
then $\forall \delta \in (0, 1)$, with probability $1-\delta$, we have $\cL(\vtheta_{i-1})-\cL^* \leq \psi\left(i, \frac{\delta}{2k}\right)$ for all $i \in [k]$, and that


$$\left|\sum_{i=1}^{k}\gamma^{k-i}X_i\right| \leq C_4\sqrt{\sum_{i=1}^{k}\gamma^{2k-2i} \psi\left(i,\frac{\delta}{2k}\right)\log\frac{4}{\delta}}$$
for some constant $C_4$.
\end{lemma}
\begin{remark}
    The $\{X_i\}$ here may not necessarily equal the $\{X_i\}$ defined in \cref{descent-direction}; we just make it general to benefit future steps. In fact, when we leverage this lemma later, we will multiply that of \cref{descent-direction} by some scalar $\in (0, 1)$.
\end{remark}  
\begin{proof}
    Note that $\sum_{i=1}^{k}\gamma^{k-i} X_i$ is a sum of martingale differences. Moreover, since $\cL$ is $\rho$-smooth and $\exists C_{1b}$ s.t. every $|X_i|$ is bounded by $C_{1b}\lnormtwo{\nabla \cL(\vtheta_{i-1})}$ (\cref{descent-direction}), we have 
    \begin{align*}
        \left|X_i\right| &\leq C_{1b}\lnormtwo{\nabla \cL(\vtheta_{i-1})} \\
        &\leq C_{1b}\sqrt{2\rho\left(\cL(\vtheta_{i-1})-\cL^*\right)} \\
        &\leq C_{1b} \sqrt{2\rho \psi(i,\delta')} \quad \text{ if }\cL(\vtheta_{i-1})-\cL^* \leq \psi(i, \delta').
    \end{align*}
    Since $\cL(\vtheta_{i-1})-\cL^* \leq \psi(i, \delta')$ holds with probability $1-\delta'$ instead of probability $1$, we create a new martingale difference sequence that masks out all the positions that exceed the bound. Specifically, we define $X'_{i,\delta'}$ as:
    $$X'_{i,\delta'} = \begin{cases}
        X_{i} & \text{ if }\cL(\vtheta_{i-1})-\cL^* \leq \psi(i, \delta'), \\
        0 & \text{ else.}
    \end{cases}$$
    This ensures that $\left|X'_{i,\delta'} \right| \leq C_{1b} \sqrt{2\rho \psi(i,\delta')}$ a.s. Then Azuma-Hoeffding's inequality gives us that for any $\epsilon'$, 
    \begin{align*}
        \mathbb{P}\left[\left|\sum_{i=1}^{k}\gamma^{k-i}X'_{i,\delta'}\right| \geq \epsilon'\right] \leq 2\exp\left(\frac{-\epsilon'^2}{4\sum_{i=1}^{k}C_{1b}^2\gamma^{2k-2i}\rho \psi(i,\delta')}\right),
    \end{align*}
    denoting the right hand side as $\frac{\delta}{2}$ gives that for any $\delta$, with probability $1-\frac{\delta}{2}$, 
    \begin{align*}
        \left|\sum_{i=1}^{k}\gamma^{k-i}X'_{i,\delta'}\right| \leq \sqrt{4\sum_{i=1}^{k}C_{1b}^2\gamma^{2k-2i}\rho \psi(i,\delta')\log\frac{4}{\delta}}.
    \end{align*}
    Let $\delta' = \frac{\delta}{2k}$, by union bound, $\cL(\vtheta_{i-1})-\cL^* \leq \psi\left(i, \frac{\delta}{2k}\right)$ for all $i \in [k]$ with probability $1-\frac{\delta}{2}$, which also implies $X'_{i,\delta'} = X_i$ for all $i \in [k]$ . So with probabilty $1-\delta$, the following two statements hold simultaneously for all $i \in [k]$:
    $$\cL(\vtheta_{i-1})-\cL^* \leq \psi\left(i, \frac{\delta}{2k}\right)$$
    and
    \begin{align*}
        \left|\sum_{i=1}^{k}\gamma^{k-i}X_i\right| & \leq \sqrt{4\sum_{i=1}^{k}C_{1b}^2\gamma^{2k-2i}\rho \psi(i,\delta')\log\frac{4}{\delta}} \\
        &= C_4\sqrt{\sum_{i=1}^{k}\gamma^{2k-2i} \psi\left(i,\frac{\delta}{2k}\right)\log\frac{4}{\delta}},
    \end{align*}
    where $C_4 = 2C_{1b}\sqrt{\rho}$.
\end{proof}


\begin{lemma}[Convergence Bound of the AGM Framework]\label{convergence}
Let $\cL$ satisfy \cref{amp:smooth}, and assume that $\mu$-PL condition is satisfied at all $\vtheta_k$ where $k \geq 0$. Let $\eta$ be a small learning rate satisfying $\frac{\beta_1}{\gamma} = \beta_1 / (1-\frac{2\eta\mu(1-\beta_1)}{R_0})\leq 0.95$. Let $K = \O(\mathrm{poly}(1/\eta))$. Under mild restrictions on $K$, for any $k \leq K$, $\delta \in (0,1)$, it holds with probability at least $1-\delta$ that $$\cL(\vtheta_{k})-\cL^* \leq C_{5a} \cdot \gamma^k \left(\cL(\vtheta_0) - \cL^*\right) + C_{5b} \cdot \eta \log \frac{K}{\delta}$$
    for some constants $C_{5a}$, $C_{5b}$.
\end{lemma}
\begin{proof}
    Denote $D_0 := \cL(\vtheta_0) - \cL^*$, and denote the bound with $1-\delta$ probability as $\psi(k, \delta) := \gamma
    ^k C_{5a}D_0 + C_{5b} \eta \log \frac{K}{\delta}$, where the constants $C_{5a}, C_{5b}$ will be specified by us later.
    We prove by induction. When $k=0$, the inequality  
$$C_{5a}D_0 + C_{5b} \eta \log \frac{K}{\delta} \geq D_0$$
holds trivially as long as $C_{5a} \geq 1$. Now assume that the statement holds for $0, 1, \cdots, k - 1$. From \cref{descent-lemma-exp}, we have
\begin{align*}
        \cL(\vtheta_{k})-\cL^* & \leq \gamma^k \left(\cL(\vtheta_0) - \cL^*\right) + \eta (1-\beta_1)\sum_{i=1}^{k}X_i  \sum_{j=i}^{k}\gamma^{k-j}\beta_1^{j-i} + C_{3} \eta.
\end{align*}
We can bound the coefficients by \begin{align*}
    \sum_{j=i}^{k}\gamma^{k-j}\beta_1^{j-i} &= \gamma^{k-i}\sum_{j=0}^{k-i}\left(\frac{\beta_1}{\gamma}\right)^j \\
    &\leq \gamma^{k-i} \cdot \frac{1}{1 - \frac{\beta_1}{\gamma}} \\
    &\leq 20 \gamma^{k-i},
\end{align*}
where the last inequality is due to the assumption $\frac{\beta_1}{\gamma} \leq 0.95$ in the statement.
Let $\tilde{X}_i := \frac{\sum_{j=i}^{k}\gamma^{k-j}\beta_1^{j-i}}{20\gamma^{k-i}}X_i$, then $\left\{\tilde{X}_i\right\}_{i=1}^{k}$ is also a martingale difference sequence and $$\left|\tilde{X}_i\right| \leq \left|X_i\right| \leq C_{1b}\lnormtwo{\nabla\cL(\vtheta_i)}\;\text{a.s.}$$
From \cref{martingale-bound-by-L}, with probability $1-\delta$, $\cL(\vtheta_{i-1})-\cL^* \leq \psi\left(i, \frac{\delta}{2k}\right)$ holds for all $i \in [k]$ and it also holds that
$$ \left|\sum_{i=1}^{k}\gamma^{k-i}\tilde{X}_i\right| \leq C_4\sqrt{\sum_{i=1}^{k}\gamma^{2k-2i} \psi\left(i,\frac{\delta}{2k}\right)\log\frac{4}{\delta}}.$$ 
The above arguments give 
\begin{align*}
    &\quad \eta (1-\beta_1)\sum_{i=1}^{k}X_i  \sum_{j=i}^{k}\gamma^{k-j}\beta_1^{j-i} \\ &\leq 20C_4\eta(1-\beta_1)\sqrt{\sum_{i=1}^{k}\gamma^{2k-2i} \psi\left(i,\frac{\delta}{2k}\right)\log\frac{4}{\delta}} \\
    &\leq 20C_4\eta(1-\beta_1)\sqrt{\sum_{i=1}^{k}\gamma^{2k-2i} \left(\gamma
    ^i C_{5a}D_0 + C_{5b} \eta \log \frac{2kK}{\delta}\right)\log\frac{4}{\delta}} \\
    &\leq 20C_4\eta(1-\beta_1)\sqrt{\sum_{i=1}^{k}\gamma^{2k-i} C_{5a}D_0 + \sum_{i=1}^{k}\gamma^{2k-2i} C_{5b}\eta\log\frac{2K^2}{\delta}}\cdot \sqrt{\log\frac{4}{\delta}} \\
    &\leq 20C_4\eta(1-\beta_1)\sqrt{\frac{\gamma^{k}C_{5a}D_0}{1-\gamma} + \frac{C_{5b}\eta}{1-\gamma^2}\log\frac{2K^2}{\delta}}\cdot \sqrt{\log\frac{4}{\delta}} \\
    &\leq 20C_4\eta(1-\beta_1)\left(\sqrt{\frac{\gamma^{k}C_{5a}D_0}{1-\gamma}} + \sqrt{\frac{C_{5b}\eta}{1-\gamma^2}\log\frac{2K^2}{\delta}} \right) \cdot \sqrt{\log\frac{4}{\delta}}.
\end{align*}
As long as $K \geq \max\{2\delta^2, 4\}$ (which is a mild restriction on $K$), we have $\log\frac{2K^2}{\delta}\log\frac{4}{\delta} \leq 3\log^2\frac{K}{\delta}$ and $\log\frac{4}{\delta} \leq \log{\frac{K}{\delta}}$. Plugging in $\frac{1}{1-\gamma} = \frac{R_0}{2\mu(1-\beta_1)}\cdot \frac{1}{\eta}$, we have
\begin{align*}
    &\quad \eta (1-\beta_1)\sum_{i=1}^{k}X_i  \sum_{j=i}^{k}\gamma^{k-j}\beta_1^{j-i} \\
    &\leq 20C_4(1-\beta_1)\left(\sqrt{\frac{C_{5a}R_0}{2\mu(1-\beta_1)}}\cdot \sqrt{\eta\gamma^{k}D_0\log\frac{K}{\delta}}+\eta\sqrt{\frac{3C_{5b}R_0}{2\mu(1-\beta_1)}\log^2\frac{K}{\delta}}\right) \\
    & \leq 10C_4 \sqrt{\frac{2C_{5a}R_0(1-\beta_1)}{\mu}}\cdot \sqrt{\eta\gamma^{k}D_0\log\frac{K}{\delta}} + 10C_4 \sqrt{\frac{6C_{5b}R_0(1-\beta_1)}{\mu}}\cdot \eta \log \frac{K}{\delta} \\
    &\leq C_{5c}\gamma^k D_0 + C_{5d}\eta\log \frac{K}{\delta},
\end{align*}
where \begin{align*}
    C_{5c} &=  5C_4\sqrt{2C_{5a}R_0(1-\beta_1)/\mu}, \\ C_{5d} &=  5C_4\sqrt{2C_{5a}R_0(1-\beta_1)/\mu} + 10C_4\sqrt{6C_{5b}R_0(1-\beta_1)/\mu}.\end{align*}Now as long as $K \geq e\delta$ (so that $\log\frac{K}{\delta}\geq 1$), we have 
\begin{align*}
     \cL(\vtheta_{k})-\cL^* & \leq \gamma^k \left(\cL(\vtheta_0) - \cL^*\right) + \eta (1-\beta_1)\sum_{i=1}^{k}X_i  \sum_{j=i}^{k}\gamma^{k-j}\beta_1^{j-i} + C_{3} \eta \\
     &\leq  \gamma^k D_0 + C_{5c}\gamma^k D_0 + C_{5d}\eta\log \frac{K}{\delta} + C_{3} \eta \\
     & \leq \left(C_{5c}+1\right)\gamma^k D_0 + \left(C_{5d}+C_3\right)\eta\log\frac{K}{\delta}.
\end{align*}
To complete the induction, we need $C_{5a}, C_{5b}$ satisfy
$$\left\{\begin{array}{lll}
    C_{5a} &\geq C_{5c}+1 &= 5C_4\sqrt{\frac{2C_{5a}D_0R_0(1-\beta_1)}{\mu}}+1,  \\
    C_{5b} &\geq C_{5d}+C_{3} &= 5C_4\sqrt{\frac{2C_{5a}D_0R_0(1-\beta_1)}{\mu}} + 10C_4\sqrt{\frac{6C_{5b}R_0(1-\beta_1)}{\mu}} + C_3. 
\end{array}\right.$$
Notice that the right-hand side grows at the rate of the square root of $C_{5a}$ and $C_{5b}$, so there must exist some feasible constants $C_{5a}$ and $C_{5b}$. Summarizing, under mild restrictions $K \geq \max\left\{2\delta^2, e\delta, 4\right\}$, the statement $\cL(\vtheta_{k})-\cL^* \leq \gamma
    ^k C_{5a}D_0 + C_{5b} \eta \log \frac{K}{\delta}$ holds with probability $1-\delta$, completing the induction.
\end{proof}


% \begin{lemma}[Convergence Bound of the AGM Framework]\label{convergence-deleted}
% Let $\cL$ satisfy \cref{amp:smooth}, and assume that $\mu$-PL condition is satisfied at all $\vtheta_k$ where $k \geq 0$. Let $\eta$ be a small learning rate satisfying $\frac{\beta_1}{\gamma} = \beta_1 / (1-\frac{2\eta\mu(1-\beta_1)}{R_0})\leq 0.95$. Let $K = \O(\mathrm{poly}(1/\eta))$. Under mild restrictions on $K$, for any $k \leq K$, $\delta \in (0,1)$, it holds with probability at least $1-\delta$ that $$\cL(\vtheta_{k})-\cL^* \leq \left(C_{5a}\gamma^k + C_{5b}\eta\right)\log \frac{K}{\delta}$$
%     for some constants $C_{5a}$ and $C_{5b}$.
% \end{lemma}
% \begin{proof}
%     We denote the bound with $1-\delta$ probability as $\psi(k, \delta) := \left(C_{5a} \gamma^k + C_{5b} \eta\right)\log \frac{K}{\delta}$, where the constants $C_{5a}, C_{5b}$ will be specified by us later.
%     We prove by induction. When $k=0$, we need  
% $$\left(C_{5a} + C_{5b} \eta\right) \left(\log K + \log \frac{1}{\delta}\right) \geq \cL(\vtheta_{0})-\cL^*, $$
% where setting $C_{5a} = \frac{\cL(\vtheta_{0})-\cL^*}{\log K}$ suffices. Now assume that the statement holds for $0, 1, \cdots, k - 1$. From \cref{descent-lemma-exp}, we have
% \begin{align*}
%         \cL(\vtheta_{k})-\cL^* & \leq \gamma^k \left(\cL(\vtheta_0) - \cL^*\right) + \eta (1-\beta_1)\sum_{i=1}^{k}X_i  \sum_{j=i}^{k}\gamma^{k-j}\beta_1^{j-i} + C_{3} \eta.
% \end{align*}
% We can bound the coefficients by \begin{align*}
%     \sum_{j=i}^{k}\gamma^{k-j}\beta_1^{j-i} &= \gamma^{k-i}\sum_{j=0}^{k-i}\left(\frac{\beta_1}{\gamma}\right)^j \\
%     &\leq \gamma^{k-i} \cdot \frac{1}{1 - \frac{\beta_1}{\gamma}} \\
%     &\leq 20 \gamma^{k-i},
% \end{align*}
% where the last inequality is due to the assumption $\frac{\beta_1}{\gamma} \leq 0.95$ in the statement.
% Let $\tilde{X}_i := \frac{\sum_{j=i}^{k}\gamma^{k-j}\beta_1^{j-i}}{20\gamma^{k-i}}X_i$, then $\left\{\tilde{X}_i\right\}_{i=1}^{k}$ is also a martingale difference sequence and $$\left|\tilde{X}_i\right| \leq \left|X_i\right| \leq C_{1b}\lnormtwo{\nabla\cL(\vtheta_i)}\;\text{a.s.}$$
% From \cref{martingale-bound-by-L}, with probability $1-\delta$, $\cL(\vtheta_{i-1})-\cL^* \leq \psi\left(i, \frac{\delta}{2k}\right)$ holds for all $i \in [k]$ and it also holds that
% $$ \left|\sum_{i=1}^{k}\gamma^{k-i}\tilde{X}_i\right| \leq C_4\sqrt{\sum_{i=1}^{k}\gamma^{2k-2i} \psi\left(i,\frac{\delta}{2k}\right)\log\frac{4}{\delta}}.$$ 
% The above arguments give 
% \begin{align*}
%     &\quad \eta (1-\beta_1)\sum_{i=1}^{k}X_i  \sum_{j=i}^{k}\gamma^{k-j}\beta_1^{j-i} \\ &\leq 20C_4\eta(1-\beta_1)\sqrt{\sum_{i=1}^{k}\gamma^{2k-2i} \psi\left(i,\frac{\delta}{2k}\right)\log\frac{4}{\delta}} \\
%     &\leq 20C_4\eta(1-\beta_1)\sqrt{\sum_{i=1}^{k}\gamma^{2k-2i} \left(C_{5a}\gamma^i + C_{5b}\eta\right)\log\frac{2kK}{\delta}\log\frac{4}{\delta}} \\
%     &\leq 20C_4\eta(1-\beta_1)\sqrt{\sum_{i=1}^{k}\gamma^{2k-2i} C_{5a}\gamma^i + \sum_{i=1}^{k}\gamma^{2k-2i} C_{5b}\eta}\cdot \sqrt{\log\frac{2K^2}{\delta}\log\frac{4}{\delta}} \\
%     &\leq 20C_4\eta(1-\beta_1)\sqrt{\frac{C_{5a}\gamma^{k}}{1-\gamma} + \frac{C_{5b}\eta}{1-\gamma^2}}\cdot \sqrt{\log\frac{2K^2}{\delta}\log\frac{4}{\delta}} \\
%     &\leq 20C_4\eta(1-\beta_1)\left(\sqrt{\frac{C_{5a}\gamma^{k}}{1-\gamma} }+ \sqrt{\frac{C_{5b}\eta}{1-\gamma}}\right)\cdot \sqrt{\log\frac{2K^2}{\delta}\log\frac{4}{\delta}}.
% \end{align*}
% As long as $K \geq \max\{2\delta^2, 4\}$ (which is a mild restriction on $K$), we have $\sqrt{\log\frac{2K^2}{\delta}\log\frac{4}{\delta}} \leq \sqrt{3\log^2\frac{K}{\delta}}$. Plugging in $\frac{1}{1-\gamma} = \frac{R_0}{2\mu(1-\beta_1)}\cdot \frac{1}{\eta}$, we have

% \begin{align*}
%     &\quad \eta (1-\beta_1)\sum_{i=1}^{k}X_i  \sum_{j=i}^{k}\gamma^{k-j}\beta_1^{j-i} \\
%     &\leq 20C_4(1-\beta_1)\left(\sqrt{\frac{C_{5a}R_0}{2\mu(1-\beta_1)}}\cdot \sqrt{\eta\gamma^{k}}+\sqrt{\frac{C_{5b}R_0}{2\mu(1-\beta_1)}}\cdot \eta\right)\cdot \sqrt{3}\log\frac{K}{\delta} \\
%     & \leq 10C_4 \sqrt{\frac{6C_{5a}R_0(1-\beta_1)}{\mu}}\cdot \sqrt{\eta\gamma^k} \log \frac{K}{\delta} + 10C_4 \sqrt{\frac{6C_{5b}R_0(1-\beta_1)}{\mu}}\cdot \eta \log \frac{K}{\delta} \\
%     &\leq \left(C_{5c}\gamma^k + C_{5d}\eta\right)\log \frac{K}{\delta},
% \end{align*}
% where \begin{align*}
%     C_{5c} &=  5C_4\sqrt{6C_{5a}R_0(1-\beta_1)/\mu}, \\ C_{5d} &=  5C_4\sqrt{6C_{5a}R_0(1-\beta_1)/\mu} + 10C_4\sqrt{6C_{5b}R_0(1-\beta_1)/\mu}.\end{align*}Now as long as $K \geq e\delta$ (so that $\log\frac{K}{\delta}\geq 1$), we have 
% \begin{align*}
%      \cL(\vtheta_{k})-\cL^* & \leq \gamma^k \left(\cL(\vtheta_0) - \cL^*\right) + \eta (1-\beta_1)\sum_{i=1}^{k}X_i  \sum_{j=i}^{k}\gamma^{k-j}\beta_1^{j-i} + C_{3} \eta \\
%      &\leq  \gamma^k \left(\cL(\vtheta_0) - \cL^*\right) + \left(C_{5c}\gamma^k + C_{5d}\eta\right)\log \frac{K}{\delta} + C_{3} \eta \\
%      & \leq \left(C_{5c}+\cL(\vtheta_0) - \cL^*\right)\gamma^k\log\frac{K}{\delta} + \left(C_{5d}+C_3\right)\eta\log\frac{K}{\delta}.
% \end{align*}
% To complete the induction, we need $C_{5a}, C_{5b}$ satisfy
% $$\left\{\begin{array}{lll}
%     C_{5a} &\geq C_{5c}+\cL(\vtheta_0) - \cL^* &= 5C_4\sqrt{\frac{6C_{5a}R_0(1-\beta_1)}{\mu}}+\cL(\vtheta_0) - \cL^*  \\
%     C_{5b} &\geq C_{5d}+C_{3} &= 5C_4\sqrt{\frac{6C_{5a}R_0(1-\beta_1)}{\mu}} + 10C_4\sqrt{\frac{6C_{5b}R_0(1-\beta_1)}{\mu}} + C_3. 
% \end{array}\right.$$
% Notice that the right-hand side grows at the rate of the square root of $C_{5a}$ and $C_{5b}$, so there must exist some feasible constants $C_{5a}$ and $C_{5b}$. Summarizing, under mild restrictions $K \geq \max\left\{2\delta^2, e\delta, 4\right\}$, the statement $\cL(\vtheta_{k})-\cL^* \leq \left(C_{5a} \gamma^k + C_{5b} \eta\right)\log \frac{K}{\delta}$ holds with probability $1-\delta$, completing the induction.
% \end{proof}

\begin{remark}\label{remark-0.95-mild}
    The assumption in the statement, $\frac{\beta_1}{\gamma} \leq 0.95$, is very mild since $\beta_1 \leq 0.9$ (\cref{amp:beta1-bound}) and $1-\gamma$ equals a constant multiple of $\eta$, so with small $\eta$ this condition is very easy to satisfy. Moreover, similar to \cref{remark-0.9-mild}, the assumed threshold $0.95$ can be replaced by any constant below $1$, and the order of the convergence rate will remain unaffected.
\end{remark}



% \xinghan{\color{orange} done complete this part.}

\subsection{Proof of Convergence-Related Conclusions in \Aref{sec:formal-statement}}\label{app:proof-formal-statement}
\begin{proof}[\textbf{Proof of \cref{thm:formal-convergence-1}}]\label{proof-formal-main-result}
This is a direct corollary following from
 \Cref{convergence}. The loss function $\cL$ is global $\mu$-PL in this case. Setting $k=K$, letting $\gamma^{K} = \O(\eta)$ gives $K = \O\left(\frac{1}{\eta}\log\frac{1}{\eta}\right)$, completing the proof.
\end{proof}

Now we move on to prove the first part of \cref{thm:formal-main-result}. The main difficulty comes from the fact that $\cL$ is only guaranteed to satisfy $\mu$-PL condition within some neighborhood $\Gamma^{\epsilon_3}$; The iteration, once getting out of that neighborhood, cannot be characterized. Hence we need to bound the probability of that event.

The trick here is to construct a proxy loss function $\tilde{\cL}$ that, agrees with $\cL$ near $\Gamma$ but has a ``wall of quadratic functions'' upon $\cL$ further away. $\tilde{\cL}$ thus satisfies $(\mu, \bar{L})$-PL which allows us to use \cref{convergence}. If the losses at all steps are small, this ensures that the iteration never leaves $\Gamma^{\epsilon_1}$, where $\cL$ and $\tilde{\cL}$ are identical. We formalize this idea in the sequel.


\begin{lemma}[Tubular Neighborhood Theorem; Theorem~6.24 in \citet{lee2012introduction}]\label{lem:tube}
Let $\Gamma\subset\R^d$ satisfy \cref{amp:low-rank-manifold} (in particular, $\Gamma$ is a $\cC^\infty$ compact embedded submanifold).
Then there exists $\tau_{\Gamma}>0$ such that, writing $N\Gamma$ for the normal bundle,
\begin{align*}
V \;:=\; \{(p,\vu)\in N\Gamma : \normtwo{\vu}<\tau_{\Gamma}\},\qquad
E:V\to\R^d,\quad E(p,\vu)=p+\vu,
\end{align*}
the map $E$ is a diffeomorphism onto the open tube $U:=\Gamma^{\tau_{\Gamma}}$.
Consequently, the nearest–point projection $P:U\to\Gamma$ is well-defined and $C^\infty$, and every $\vtheta\in U$ can be written uniquely as
\begin{align*}
\vtheta \;=\; P(\vtheta)+\nu(\vtheta), \quad \nu(\vtheta)\in N_{P(\vtheta)}\Gamma.
\end{align*}
\end{lemma}

\begin{corollary}[Smooth distance and unit normal on the tube]\label{cor:dist-normal}
In the setting of \cref{lem:tube}, let $U:=\Gamma^{\tau_\Gamma}$ and write
$E^{-1}(\vtheta)=(P(\vtheta),\nu(\vtheta))$ on $U$. Define
\begin{align*}
  r(\vtheta):=\dist(\vtheta,\Gamma)=\normtwo{\nu(\vtheta)},\qquad
  \vn(\vtheta):=\frac{\nu(\vtheta)}{\normtwo{\nu(\vtheta)}} \quad (\vtheta\in U\setminus\Gamma).
\end{align*}
Then $r$ and $\vn$ are $\cC^\infty$  on $U\setminus\Gamma$, and
\begin{align}\label{eq:grad-r-equals-n}
  \nabla r(\vtheta)=\vn(\vtheta), \qquad\forall\vtheta\in U\setminus\Gamma.
\end{align}
\end{corollary}

\begin{proof}By \cref{lem:tube}, $P$ and $\nu$ are $\cC^\infty$ on $U$, hence so are $r$ and $\vn$ on $U\setminus\Gamma$.
Set $g(\vtheta):=r(\vtheta)^2=\normtwo{\nu(\vtheta)}^2$. 
By the chain rule,
\begin{align*}
  \nabla g(\vtheta)
  \;=\; 2\,\big(\nabla \nu(\vtheta)\big)^\top \nu(\vtheta).
\end{align*}
From the identity $\vtheta=P(\vtheta)+\nu(\vtheta)$ we have \begin{align*}
  \nabla \nu(\vtheta)\;=\; I - \nabla P(\vtheta).
\end{align*}
Since $\nabla P(\vtheta)$ maps into $T_{P(\vtheta)}\Gamma$ and 
$\nu(\vtheta)\in N_{P(\vtheta)}\Gamma$, it follows that
\begin{align*}
  \big(\nabla P(\vtheta)\big)^\top \nu(\vtheta) \;=\; \vzero,
\end{align*}
hence
\begin{align*}
  \nabla g(\vtheta)
  \;=\; 2\,(I-\nabla P(\vtheta))^\top \nu(\vtheta)
  \;=\; 2\,\nu(\vtheta).
\end{align*}
Therefore, for $\vtheta\in U\setminus\Gamma$ (so $r(\vtheta)>0$),
\begin{align*}
  \nabla r(\vtheta)
  \;=\; \frac{1}{2\,r(\vtheta)}\,\nabla g(\vtheta)
  \;=\; \frac{\nu(\vtheta)}{\normtwo{\nu(\vtheta)}} 
  \;=\; \vn(\vtheta),
\end{align*}
which is \Cref{eq:grad-r-equals-n}.
\end{proof}

\begin{lemma}[Nonobtuse angle between $\nabla\cL$ and the outward normal]\label{lem:acute-angle-short}
Let $\Gamma$ satisfy \cref{amp:low-rank-manifold} and let $\cL$ satisfy \cref{amp:basic} and \cref{amp:smooth}.
Then there exists a constant $\tau\in(0,\tau_{\Gamma}]$ such that, for all
$\vtheta\in\Gamma^{\tau}$ with nearest-point projection $P(\vtheta)$, distance
$r(\vtheta):=\normtwo{\vtheta-P(\vtheta)}$, and outward unit normal
$\vn(\vtheta):=(\vtheta-P(\vtheta))/r(\vtheta)$, we have
\begin{align}\label{eq:acute-angle-claim}
\big\langle \nabla \cL(\vtheta),\, \vn(\vtheta)\big\rangle \geq 0.
\end{align}
In other words, $\angle\!\big(\nabla\cL(\vtheta),\vn(\vtheta)\big)\leq\pi/2$ on
$\Gamma^{\tau}\setminus\Gamma$.
\end{lemma}

\begin{remark}
    This lemma also implies that $t\mapsto \cL\big(P(\vtheta)+t\vn(\vtheta)\big)$ is non-decreasing on $(0,\tau]$, since $\frac{d}{dt}\cL\big(\vphi+t\vn\big)=\langle \nabla\cL(\vphi+t\vn),\vn\rangle \ge 0$ on $(0,\tau]$. To put it vividly, $\Gamma^{\tau}$ is a valley, with $\Gamma$ being the floor at the center of it.
\end{remark}
\begin{proof}
By \cref{amp:low-rank-manifold}, each $\vzeta\in\Gamma$ is a local minimizer of $\cL$, hence $\nabla\cL(\vzeta)=\vzero$,
and $\nabla^2\cL(\vzeta)$ is positive definite on $N_{\vzeta}\Gamma$, with a uniform lower-bound
$m>0$ of its eigenvalues on $N_{\vzeta}\Gamma$ (from the compactness of $\Gamma$). Let $\tau_{\Gamma}$ be as in \cref{lem:tube}. 
For $\vtheta\in \Gamma^{\tau_{\Gamma}}$ write $\vphi:=P(\vtheta)$, $r:=\normtwo{\vtheta-\vphi}$,
and $\vn:=(\vtheta-\vphi)/r\in N_{\vphi}\Gamma$. Since $\cL$ is $\cC^5$-smooth, we can perform a third order Taylor expansion:
\begin{align*}
\nabla\cL(\vtheta)
= \nabla\cL(\vphi) + \nabla^2\cL(\vphi)(\vtheta-\vphi) + \vtheta^{(3)},
\qquad \normtwo{\vtheta^{(3)}} \leq C^{(3)} r^2,
\end{align*}
where $C^{(3)}$ is a constant independent of $\vtheta$.
Taking the inner product with $\vn$ and using $\nabla\cL(\vphi)=\vzero$ and $\vtheta-\vphi= r\vn$, we have
\begin{align*}
\big\langle \nabla\cL(\vtheta), \vn\big\rangle
= r\vn^\top\nabla^2\cL(\vphi)\vn + \big\langle \vtheta^{(3)},\vn\big\rangle
\geq m r - C^{(3)} r^2 .
\end{align*}
Choose $\tau:=\min\{\tau_{\Gamma},\, m/C^{(3)}\}$, 
then for all $r\le\tau$, we obtain \Cref{eq:acute-angle-claim}.
\end{proof}

\begin{proof}[\textbf{Proof of the first part of \cref{thm:formal-main-result}}] First we construct a tubular neighborhood around $\Gamma$ and introduce some notations.
By \cref{amp:low-rank-manifold} , $\Gamma$ is the unique set of minimizers of $\cL$ in some  neighborhood $U$ of $\Gamma$. 
By \Cref{lem:acute-angle-short}, there is a $\tau\in(0,\tau_\Gamma]$ for which the
nonobtuse condition $\langle \nabla\cL(\vtheta),\vn(\vtheta)\rangle\ge 0$ holds on
$\Gamma^\tau\setminus\Gamma$. By \Cref{lem:working-zone}, there exists
$\epsilon_3>0$ such that $\cL$ is $\mu$-PL on $\Gamma^{\epsilon_3}$. Shrinking $\epsilon_3$ if necessary, assume $\epsilon_3 \leq \tau \text{ and } \Gamma^{\epsilon_3} \subseteq U$.

Throughout this proof we work inside the tube $\Gamma^{\tau_\Gamma}$ given by
\Cref{lem:tube}. In particular, for every $\vtheta\in \Gamma^{\tau_\Gamma}$ we have the nearest–point
projection $P(\vtheta)\in\Gamma$, the normal offset $\nu(\vtheta)\in N_{P(\vtheta)}\Gamma$,
the distance and unit normal
\begin{align*}
  r(\vtheta):=\dist(\vtheta,\Gamma)=\normtwo{\nu(\vtheta)},\qquad
  \vn(\vtheta):=\frac{\nu(\vtheta)}{\normtwo{\nu(\vtheta)}},
\end{align*}
and (on $U\setminus\Gamma$) the identity $\nabla r(\vtheta)=\vn(\vtheta)$ from
\Cref{cor:dist-normal}.

Next, recall the constant $\epsilon_1$ constructed in \cref{lem:working-zone} such that $0<\epsilon_1<\epsilon_3$. 
Define the “gap level”
\begin{align*}
  \cL_m := \min\left\{\inf\big\{ \cL(\vtheta):r(\vtheta)\in[\epsilon_1,\epsilon_3] \big\}, \bar{L} \right\}.
\end{align*}
The set $\{\vtheta:\ r(\vtheta)\in[\epsilon_1,\epsilon_3]\}$ is compact and
disjoint from $\Gamma$, hence $\cL_m>\cL^*$.

Then we define the proxy objective $\tilde{\cL}:\R^d\to\R$ by
\begin{align*}
  \tilde{\cL}(\vtheta)
  :=\begin{cases}
      \cL(\vtheta), & \dist(\vtheta,\Gamma)\le \epsilon_1,\\[2pt]
      \cL(\vtheta)+\dfrac{C}{2}\big(\dist(\vtheta,\Gamma)-\epsilon_1\big)^2, & \dist(\vtheta,\Gamma)>\epsilon_1,
    \end{cases}
\end{align*}
where $C$ is a large constant satisfying $C \geq \mu$. Note that for $\vtheta \in \R^d$,
\[
  \tilde{\cL}(\vtheta) < \cL_m \;\Longrightarrow\; \mathrm{dist}(\vtheta)<\epsilon_1,
\]
so on the sublevel set $\{\tilde{\cL}<\cL_m\}$ we have $\tilde{\cL}=\cL$.

Now define \begin{align*}
    \bar{L} := \frac{C}{2}(\epsilon_3-\epsilon_1)^2 + \cL^*, 
\end{align*}and we prove the core property of the proxy loss function: $\tilde{\cL}$ is $(\mu,\bar{L})$-PL. First we consider the case $\vtheta \in \Gamma^{\epsilon_3}$. On $\Gamma^{\epsilon_3}$, the distance function $r(\vtheta) = \dist(\vtheta,\Gamma)$ is defined. Using $\nabla r(\vtheta)=\vn(\vtheta)$ from \Cref{cor:dist-normal} and the
nonobtuse condition from \Cref{lem:acute-angle-short}, for $r(\vtheta)>\epsilon_1$,
\begin{align*}
  \nabla \tilde{\cL}(\vtheta)
  &= \nabla \cL(\vtheta) + C\big(r(\vtheta)-\epsilon_1\big)\,\vn(\vtheta),\\
  \|\nabla \tilde{\cL}(\vtheta)\|_2^2
  &= \|\nabla\cL(\vtheta)\|_2^2
     + 2C\big(r(\vtheta)-\epsilon_1\big)\,\langle \nabla\cL(\vtheta),\vn(\vtheta)\rangle
     + C^2\big(r(\vtheta)-\epsilon_1\big)^2 \\
  &\ge \|\nabla\cL(\vtheta)\|_2^2 + C^2\big(r(\vtheta)-\epsilon_1\big)^2.
\end{align*}
Since $\cL$ is $\mu$-PL on $\Gamma^{\epsilon_3}$,
\begin{align*}
  \|\nabla \tilde{\cL}(\vtheta)\|_2^2
  \;\ge\; 2\mu\big(\cL(\vtheta)-\cL^*\big) + 2C\cdot \frac{C}{2}\big(r(\vtheta)-\epsilon_1\big)^2
  \;\ge\; 2\min\{\mu,C\}\big(\tilde{\cL}(\vtheta)-\cL^*\big).
\end{align*}
For $r(\vtheta)\le \epsilon_1$, $\tilde{\cL}=\cL$ and the $\mu$-PL inequality holds
trivially. Our choice of $C$ yields
\begin{align*}
  \|\nabla \tilde{\cL}(\vtheta)\|_2^2
  \;\ge\; 2\mu\big(\tilde{\cL}(\vtheta)-\cL^*\big)
  \qquad \text{for all }\ \vtheta\in\Gamma^{\epsilon_3},
\end{align*}
i.e., $\tilde{\cL}$ is $\mu$-PL on $\Gamma^{\epsilon_3}$. 

Combining with the fact that $\tilde{\cL}(\vtheta) > \bar{L}$ if $\vtheta \notin \Gamma^{\epsilon_3}$, we conclude that $\tilde{\cL}$ is $(\mu,\bar{L}$)-PL. 

Finally we come back to prove the main conclusion. Let $\bar{C}$ be the constant constructed in \cref{cor:nablaL-large}. Note that for $\vtheta \in \R^d \setminus \Gamma^{\epsilon_1}$, \begin{align*}
    \normtwo{\nabla\tilde{\cL}(\vtheta)} \leq \bar{C} &\Rightarrow \normtwo{\nabla \cL(\vtheta) + C\big(r(\vtheta)-\epsilon_1\big)\,\vn(\vtheta)} \leq \bar{C} \\
     &\Rightarrow \normtwo{C\big(r(\vtheta)-\epsilon_1\big)\,\vn(\vtheta)} \leq \bar{C} \quad \text{(from \cref{lem:acute-angle-short})}\\
    &\Rightarrow r(\vtheta) \leq \frac{\bar{C}}{C} + \epsilon_1. 
\end{align*}
Increasing $C$ if necessary, we can let $\bar{C} / C < (\epsilon_3 - \epsilon_1) / 2$, which implies \begin{align*}
    \epsilon_{m} :=  \frac{\bar{C}}{C} + \epsilon_1 < \frac{\epsilon_1 + \epsilon_3}{2}.
\end{align*} Take some $\cL_0 > \cL^*$ such that $\cL_0 - \cL^* < (\cL_m - \cL^*) / 2C_{5a}$, where $C_{5a}$ is the constant in \cref{convergence}. Take a constant $\epsilon_0 < \epsilon_1$ such that all loss values inside $\Gamma^\epsilon_0$ are bounded by $\cL_0$. There exists a sufficiently small learning rate $\eta$ such that (let $K=\lfloor (T+1)\eta^{-2}\rfloor$):
\begin{enumerate}
    \item $C_{5a}\gamma^k \left(\cL_0 - \cL^*\right) + C_{5b}\eta\log \frac{K}{\delta} \leq 0.99\cL_m - \cL^*$, $\forall k \leq K, \delta \in (\eta^{200}, 1)$. 
    \item $\eta \cdot R \cdot R_2 < \epsilon_3 - \epsilon_{m}$.
\end{enumerate}
The second property ensures that any single step of update cannot jump from the interior of $\Gamma^{\epsilon_m}$ to the exterior of $\Gamma^{\epsilon_3}$. However when $\vtheta_{k-1} \in \Gamma^{\epsilon_3} \setminus \Gamma^{\epsilon_m}$, it follows that $\normtwo{\nabla\tilde{\cL}(\vtheta)} > \bar{C}$, so $\cL(\vtheta_k) < \cL(\vtheta_{k-1})$ from \cref{cor:nablaL-large}. By induction, we conclude that for any $\vtheta_0 \in \Gamma^{\epsilon_0}$, if we launch an AGM from $\vtheta_0$ and train using $\tilde{\cL}$ and $\eta$, all loss values $\tilde{\cL}(\vtheta_k), k \in [0, K-1]$ do not exceed $\bar{L}$, which means the $\mu$-PL condition of $\tilde{\cL}$ is satisfied at all $\vtheta_k$ where $k \in [0, K-1]$. This meets the requirement of \cref{convergence}.

By \cref{convergence} and the first property of $\eta$, we further conclude that all loss values $\tilde{\cL}(\vtheta_k)$ do not exceed $0.99\cL_m$. Finally, by noting that \begin{align*}
\tilde{\cL}(\vtheta) < \cL_m \Rightarrow \vtheta \in \Gamma^{\epsilon_1} \text{  and } \tilde{\cL}(\vtheta) = \cL(\vtheta),  
\end{align*}
we conclude from \cref{convergence} that: For any $\vtheta_0 \in \Gamma^{\epsilon_0}$, if we launch an AGM from $\vtheta_0$ and train using $\cL$ and $\eta$, for any $k \leq K$, $\delta \in (\eta^{200},1)$, it holds almost surely that $\vtheta_k \in \Gamma^{\epsilon_1}$, and it holds with probability at least $1-\delta$ that $$\cL(\vtheta_{k})-\cL^* \leq C_{5a} \cdot \gamma^k \left(\cL(\vtheta_0) - \cL^*\right) + C_{5b} \cdot \eta \log \frac{K}{\delta},$$
and it takes $K_0 = \mathcal{O}(\frac{1}{\eta} \log \frac{1}{\eta})$ time to reach $\cL(K_0) - \cL^* = \O\left(\eta \log \frac{1}{\eta\delta}\right)$, 
completing the proof.
% Plugging in $K=\lfloor (T+1)\eta^{-2}\rfloor$, 
% we conclude that there exists some $\epsilon < 0.5\epsilon_2$ such that if $\cL(\vtheta_{0})-\cL^* \leq \sup\left\{\cL(\vtheta) : \vtheta \in \Gamma^{\epsilon}\right\} -\cL^*$ and $\eta$ is sufficiently small, then with probability $1-\delta$, the loss values at all steps is strictly smaller than $\cL_m$, and $\eta R/\epsilon < 0.5\epsilon_2$ so any single step of update cannot jump from the interior of $\Gamma^{0.5\epsilon_2}$ to the exterior of $\Gamma^{\epsilon_2}$. So if the statement in \cref{convergence} holds, all iterations of AGM stay inside $\Gamma^{\epsilon_2}$. Then substituting $K_0 = \lceil\log_{\gamma} \eta \rceil= \O\left(\frac{1}{\eta}\log \frac{1}{\eta}\right)$ gives the result. 
\end{proof}


